\cotefile{62}
\scanpage{83}
% L1, top
\[
\sigma(u+2\omega_i)=\sigma u\,e^{2\eta_i(u+\omega_i)+2i\pi}
\tag*{\textcircled{\scriptsize 16}}
\]

% L2–4, upper quarter
\uline{Formule de Jacobi} \quad $f$ fonction elliptique. $(a_i)$ et $(b_j)$ $(i=1,\ldots p)$ système \unsure{fondamental} de zéros et de \unsure{poles|pôles}, tel que $\sum a_i=\sum b_i$. Alors

% L5
\[
f(\unsure{z|u})=c^{\mathrm{te}}\times
\frac{\prod\sigma(u-a_i)}{\prod\sigma(u-b_i)}
\tag*{\textcircled{\scriptsize 17}}
\]

% L6–7
(le deuxième mb sera en effet elliptique, \unsure{grace|grâce} à $(16)$, et à $\sum a_i=\sum b_i$)

% L8, boxed heading
5) \fbox{Formules d'addition des fonctions elliptiques}

% L9–10
La formule de Jacobi, appliquée à la fonction de $u$ \ $\unsure{p|\wp}u-pv$, donne facilement

% L11
\[
pu-pv=-\frac{\sigma(u+v)\sigma(u-v)}{\sigma^2u\,\sigma^{\unsure{2|7}}v}
\tag*{\textcircled{\scriptsize 18}}
\]

% L12–13
Prenant la dérivée logarithmique en $u$, en $v$ et en \unsure{ajoutant}, il vient

% L14–15, left brace
\[
\left\{\begin{aligned}
\frac{p'u}{pu-pv}&=\zeta(u+v)+\zeta(u-v)-2\zeta u\\
-\frac{p'v}{pu-pv}&=\zeta(u+v)-\zeta(u-v)-2\zeta v
\end{aligned}\right.
\tag*{\textcircled{\scriptsize 19}}
\]

% L16
\[
\zeta(u+v)-\zeta u-\zeta v=\frac12\,\frac{p'u-p'v}{pu-pv}
\tag*{\textcircled{\scriptsize 20}}
\]

% L17
faisant $v=u$ :

% L18
\[
\zeta(2u)-2\zeta u=\frac12\,\frac{p''u}{p'u}
\tag*{(20 bis)}
\]

% L19, foot
Autres \unsure{formule|formules}

% lower left
% archive sheet number: 82
\scanpage{84}
% L1
\[
\unsure{p|\wp}(u+v)+pu+pv=\frac14\left(\frac{p'u-p'v}{pu-pv}\right)^2
\tag*{(21)}
\]

% L2–3
(en dérivant en $u$ et $v$ l'expression de $\frac{p'u}{pu-pv}$, et y remplaçant $p''$ en fonction de $p$)

% L4
faisant $v=u$

% L5
\[
p(2u)+2pu=\frac14\left(\frac{p''u}{p'u}\right)^2
\tag*{(21 bis)}
\]

% L6
L'application de la formule d'Hermite à $p(2u)$ donne

% L7
\[
4p(2u)=pu+p(u+\omega_1)+p(u+\omega_2)+p(u+\omega_3)
\tag*{(22)}
\]

% L8, boxed heading
5) \fbox{Fonctions $pu-e_i$}

% L9–10
L'application de la formule d'Hermite à $pu-e_i$ donne

% L11–12
\[
\frac{1}{pu-e_i}=\frac{p(u+\omega_i)-e_i}{(e_j-e_i)(e_k-e_i)}
\quad(i,j,k\ \text{\unsure{permutation} de }1,2,3)
\tag*{\textcircled{\scriptsize 23}}
\]

% L13
(compte tenu de $e_1+e_2+e_3=0$)

% L14–15
La racine carrée de \deleted{\unsure{cette|cet}} fonction est analytique méromorphe; la formule d'addition $pu-pv\inserted{\unsure{\ldots|=\ldots}}$ donne

% L16
\[
pu-e_i=-\frac{\sigma(u+\omega_i)\,\sigma(u-\omega_i)}{\sigma^2u\,\sigma^2\omega_i}
=\left(\frac{\sigma_i u}{\sigma u}\right)^{\unsure{2|7}}\quad\text{où}
\tag*{\textcircled{\scriptsize 24}}
\]

% L17
\[
\sigma_i u=\frac{1}{\sigma\omega_i}\,\sigma(u+\omega_i)e^{-\eta_i u}
\tag*{\textcircled{\scriptsize 25}}
\]

% L18–21, foot
On trouve que $\sqrt{pu-e_i}=\frac{\sigma_i u}{\sigma u}$ a les périodes $2\omega_i$ et $4\omega_j$ (mais non $2\omega_i$ et $2\omega_j$) -- $i\ne j$; elle change de signe si on ajoute $2\omega_i$ à $u$. En d'autres termes

% lower left
% archive sheet number: 83
\scanpage{85}
% L1–2
\[
\left\{\begin{aligned}
\sigma_i(u+2\omega_i)&=\sigma_i u\,e^{2\eta_i(u+\omega_i)+i\pi}\\
\sigma_i(u+2\omega_j)&=\sigma_i u\,e^{2\eta_j(u+\omega_j)}
\end{aligned}\right.\qquad(j\ne i)
\tag*{(26)}
\]

% L3–5
Ainsi $\sqrt{\unsure{p|\wp}u-e_1}$ a $2\omega_1,4\omega_2$ comme système fondamental de périodes, $\sqrt{pu-e_2}$ admet $2\omega_2,4\omega_3$, et $\sqrt{pu-e_3}$ admet $2\omega_3,4\omega_1$

% L6
Les fonctions $\sigma_i u$ sont paires, $\sigma_i(\unsure{0|o})=1$.

% L7, below double rule
\par\noindent\rule{\linewidth}{.4pt}\par\vspace{-1ex}\noindent\rule{\linewidth}{.4pt}\par
\begin{center}\fbox{Les fonctions $\Theta$}\end{center}

% L8–10
S'introduisent pour le calcul numérique des fonctions elliptiques; introduisent une \unsure{dissymétrie} entre $\omega_1$ et $\omega_2$.

% L11–15
Par la substitution $u\to u+2\omega_1$, $\sigma u$ et $\deleted{e^{\unsure{2\eta_1(u+\omega_1)}}}e^{\frac{\eta_1 u^2}{2\omega_1}-\frac{i\pi u}{2\omega_1}}$ sont multipliés par le même facteur $e^{2\eta_1(u+\omega_1)-i\pi}$, donc leur quotient admet la période $2\omega_1$, et se développe donc en série de Fourier; par suite

% L16
\[
\sigma u=e^{\frac{\eta_1u^2}{2\omega_1}}\sum_{-\infty}^{+\infty} A_n\,e^{(2n-1)i\pi\frac{u}{2\omega_1}}
\tag*{1}
\]

% L17–18
En exprimant que $\sigma(u+2\omega_2)=\sigma u\,e^{2\eta_2(u+\omega_2)-i\pi}$ on obtient une relation de récurrence

% L19
\[
A_{n+1}=-{\unsure{h|b}}^nA_n,\quad\text{où}\quad h=e^{i\pi\frac{\omega_2}{\omega_1}}
\qquad\left(\unsure{\mathfrak I|J}\,\frac{\omega_2}{\omega_1}>0\right),
\]

% L20–22, foot
ce qui détermine les $A_n$ \deleted{\unsure{à|a}} au facteur $A_0$ près. L'expression de $\sigma u$ amène à introduire la fonction

% lower left
% archive sheet number: 84
\scanpage{86}
% L1
\[
\Theta(v\mid t)=-\frac{1}{\unsure{i|1}}\sum_{-\infty}^{+\infty}(-1)^n e^{(n-\frac12)^2i\pi t+(2n-1)i\pi v}
\tag*{\textcircled{\scriptsize 1}}
\]

% L2–4
série qui peut aussi s'écrire, en \deleted{\unsure{groupant}} décomposant la somme en $\sum_1^\infty$ et $\sum_{-\infty}^{\unsure{0|o}}$ et groupant les termes deux à deux

% L5
\[
\Theta(v\mid t)=2_{\deleted{\unsure{i|1}}}\sum_1^\infty(-1)^{n-1}e^{(n-\frac{1}{\unsure{1|2}})^2i\pi t}\sin(2n-1)\pi v
\tag*{\textcircled{\scriptsize 2}}
\]

% L6–7
On a alors, par construction. $\Theta$ est définie \inserted{et holomorphe} pour $\unsure{\mathfrak I|J}t>0$ et pour $v$ quelconque

% L8
\[
\sigma(u\mid\omega_1,\omega_2)=
\frac{2\omega_1}{\Theta'_v(0\mid\frac{\omega_2}{\omega_1})}\,
e^{\frac{\eta_1u^2}{2\omega_1}}\,\Theta\left(\frac{u}{2\omega_1},\frac{\omega_2}{\omega_1}\right)
\tag*{\textcircled{\scriptsize 3}}
\]

% L9–11
La fonction $\Theta$ est exprimée par une série très rapidement convergente. -- $\Theta(v\mid t)$ \unsure{est} fonction impaire de $v$, et on a

% L12–13
\[
\left\{\begin{aligned}
\Theta(v+1\mid t)&=-\Theta(v\mid t)&\Theta(v+2\mid t)&=\Theta(v\mid t)\\
\Theta(v+t\mid t)&=\Theta(v\mid t)e^{-2i\pi v-i\pi t-i\pi}
\end{aligned}\right.
\tag*{\textcircled{\scriptsize 4}}
\]

% L14–16
Ces relations caractérisent $\Theta$ \unsure{a|à} un facteur constant près, pour $t$ donné (comme on voit par le calcul qui avait donné $\Theta$)

% L17
Notons aussi la formule (vérification immédiate)

% L18
\[
\Theta(v\mid t+4)=-\Theta(v\mid t),\qquad\text{d'où}\quad\Theta(v\mid t+8)=\Theta(v\mid t)
\tag*{\textcircled{\scriptsize 5}}
\]

% L19–20
Les zéros de $\Theta(v\mid t)$ sont déterminés par ceux de $\sigma u$, ce sont les nombres $v=m_1+m_2 t$. --

% L21–22, foot
Pour une expression de \textcircled{\scriptsize 1}, \textcircled{\scriptsize 2}, commode au calcul, on pose parfois $h(t)=h=e^{i\pi t}$; alors

% lower left
% archive sheet number: 85
\scanpage{87}
% L1–3
\[
\begin{aligned}
&\left\{\begin{aligned}
\Theta(v\mid t)&=-\frac1{\unsure{i|1}}\sum_{-\infty}^{\infty}(-1)^n h(t)^{(n-\frac12)^2}e^{(2n-1)i\pi v}\\
&=2\sum_1^\infty(-1)^{n-1}h(t)^{(n-\frac12)^2}\sin(2n-1)\pi v
\end{aligned}\right.\\
&\hspace{5em}=\deleted{\unsure{2\sum|2h\sum}}\,2\left[h(t)^{\frac14}\sin\pi v-h(t)^{\frac94}\sin3\pi v+\ldots\right]
\end{aligned}
\tag*{\textcircled{\scriptsize 6}}
\]

% L4
On peut poser $t=iT$ \unsure{--|où} $\unsure{\Re|\mathcal R}T>0$, et on a

% L5
\[
\Theta(v\mid iT)=2\sum_{\deleted{\unsure{-\infty}}\unsure{1}}^{+\infty}(-1)^{n-1}e^{-(n-\frac12)^2\pi T\deleted{\unsure{+2|+i}}}\sin(2n-1)\pi v
\]

% L6–9
La fonction $\Theta(v\mid t)$ (et les séries analogues qui vont suivre) permet d'exprimer la somme d'une série de la forme $\sum_{-\infty}^{+\infty}e^{-an^2+bn+c}$, où $\Re a>0$, -- et réciproquement \unsure{bien sur|bien sûr}. On a

% L10
\[
\sum_{-\infty}^{+\infty}e^{-an^2+bn+c}=-e^{c+\frac b2-\frac a4}\,\Theta\left(\frac{b-a-i\pi}{2i\pi}\,\middle|\,\frac{ia}{\pi}\right)
\tag*{\textcircled{\scriptsize 7}}
\]

% L11, underlined heading
\uline{Formule de transformation pour $t\to-\frac1t$}

% L12–13
En exprimant que $\sigma(u\mid\omega_1,\omega_2)=\sigma(u\mid\omega_2,-\omega_1)$, on obtient une équation fonctionnelle pour $\Theta$

% L14
\[
\Theta(v\mid t)=e^{-\frac{i\pi v^2}{t}}\,\Theta\left(\frac vt\,\middle|\,-\frac1t\right)\,t\,
\frac{\Theta'_v(\unsure{0|v}\mid t)}{\Theta'_v(0\mid-\frac1t)}
\tag*{\textcircled{\scriptsize 8}}
\]

% L15–17
Pour expliciter $t\frac{\Theta'_v(0\mid t)}{\Theta'_v(0\mid-\frac1t)}=H(t)$, on note d'abord que \uline{$\Theta(v\mid t)$ est solution de l'équation de la chaleur :}

% L18
\[
\frac{\partial^{\unsure{2}}\Theta}{\partial v^2}=4i\pi\frac{\partial\Theta}{\partial t}
\tag*{\textcircled{\scriptsize 9}}
\]

% lower left
% archive sheet number: 86
\scanpage{88}
% L1–7
(vérification immédiate par la série). En appliquant au deuxième membre de $(8)$, où $t\frac{{\Theta}'_v(0\mid t)}{\Theta'_v(0\mid-\frac1t)}$ est laissé sous forme indéterminée, on trouve une relation sur $H(t)$ : $\frac{H'(t)}{H(t)}=-\frac1{2t}$, d'où suit que $H(t)$ est proportionnel à $\frac1{\sqrt t}$; le facteur de proportionnalité s'obtient en faisant dans

% L7
\[
\Theta(v\mid t)=e^{-\frac{i\pi v^2}{t}}\Theta\left(\frac vt\,\middle|\,-\frac1t\right)\frac a{\sqrt t},
\]

% L8–9
$t=i$ (d'où $t=-\frac1t$) et $v=\unsure{0|o}$, on trouve $a=\sqrt i$, d'où enfin

% L10, boxed
\[
\boxed{\Theta(v\mid t)=e^{-\frac{i\pi v^2}{t}}\Theta\left(\frac vt\,\middle|\,-\frac1t\right)i\sqrt{\frac it}}
\tag*{\textcircled{\scriptsize 10}}
\]

% L11–13
(où la détermination de $\sqrt{\frac it}$ est celle qui correspond à $-\frac{\pi}{2}\leq\operatorname{Arg}\frac it\leq+\frac{\pi}{2}$). Comparant avec \textcircled{\scriptsize 8}, on trouve donc

% L14
\[
\Theta'_v(0\mid t)=\frac{i}t\sqrt{\frac it}\,\Theta'_v(0\mid-\frac1t)
\tag*{\textcircled{\scriptsize 11}}
\]

% L15–16
Qu'on exprime généralement, faisant $t=iT$. (\unsure{où} $\unsure{\Re|\mathcal R}T>0$)

% L17
\[
\sum_1^\infty(-1)^{n-1}(2n-1)e^{-(n-\frac12)^2\pi T}
=\frac1{T^{3/2}}\sum_{\unsure{1}}^\infty(\unsure{-1})^{n-1}(2n-1)e^{-(n-\frac1{\unsure{2|3}})\frac{\pi}{\unsure{T|1}}}
\tag*{\textcircled{\scriptsize 12}}
\]

% L18
\[
\text{(car}\qquad\Theta'_v(\unsure{v|0}\mid t)=2\pi\sum_1^\infty(-1)^{n-1}(2n-1)e^{(n-\frac12)^2i\pi t}\cos(2n-1)\unsure{i\pi v|\pi v}
\tag*{\textcircled{\scriptsize \unsure{12|17} bis}}
\]

% L19, foot
\[
\text{d'où}\qquad\Theta'_v(0\mid iT)=2\pi\sum_1^\infty(\unsure{-1|1})^{n-1}(2n-1)e^{-(n-\frac12)^2\pi T}
\]

% lower left
% archive sheet number: 87
\scanpage{89}
% L1–5
La formule \textcircled{\scriptsize 10} est \unsure{importante} pour le calcul numérique des fonctions $\Theta$, car ramène le calcul de $\Theta(v\mid iT)$ à celui de $\Theta(v\mid\frac iT)$; si donc $\unsure{\Re|\mathcal R}T\ngtr1$, alors on a $\Re T>1$ ou $\Re\frac1T>1$, d'où résulte que dans l'expression

% L6
\[
\Theta(v\mid iT)=2\sum_1^\infty(-1)^{n-1}h^{(n-\frac1{\unsure{2|3}})^2}\sin(2n-1)\pi v,\quad\text{où}
\]

% L7–11
$h=e^{-\pi T}$ ou $e^{-\pi\frac1T}$, on peut supposer que $|h|\leq e^{-\pi}\mathrel{\#}\frac1{10}$; comme l'exposant de $h$ augmente rapidement, on voit qu'on a bien une série très rapidement convergente.

% L12, underlined
\uline{Autres fonctions $\Theta$ liées aux fonctions elliptiques}

% L13
La fonction $\Theta_i$ $(i=1,2,3)$ est définie par

% L14
\[
\sigma_i u=e^{\frac{\eta_1u^2}{2\omega_1}}\,
\frac{\Theta_i\left(u/2\omega_1\,\middle|\,\frac{\omega_2}{\omega_1}\right)}
{\Theta_i\left(0\,\middle|\,\frac{\omega_2}{\omega_1}\right)}
\tag*{\textcircled{\scriptsize 13}}
\]

% L15
et s'exprime par une série du type $\Theta$ :

% L16–19, brace continuing
\[
\left\{\begin{aligned}
\Theta_1(v\mid t)&=\deleted{\unsure{2|e}}\sum_{-\infty}^\infty e^{(n-\frac1{\unsure{1|2}})^2i\pi t+(2n-1)i\pi v}\\
&=2\sum_1^\infty e^{(n-\frac11)^2i\pi t}\cos(2n-1)\pi v=\vartheta_2(v)\\
\Theta_2(v\mid t)&=\sum_{-\infty}^{+\infty}(\unsure{-1|-i})^n e^{n^2i\pi t+2ni\pi v}\deleted{\unsure{\sum^{\infty}|2\sum}}\\
&=1+2\sum_1^\infty e^{n^2i\pi t}\cos2\unsure{\pi|n}\pi v=\vartheta_0(v)
\end{aligned}\right.
\tag*{\textcircled{\scriptsize 14}}
\]

% lower left
% archive sheet number: 88
\scanpage{90}
% L1–2, continuation of brace
\[
\left\{\begin{aligned}
\Theta_3(v\mid t)&=\sum_{-\infty}^{+\infty}e^{n^2i\pi t+2ni\pi v}\\
&=1+2\sum_1^\infty e^{n^2i\pi t}\cos2n\pi v=\vartheta_3(v)
\end{aligned}\right.
\]

% L3
($\Theta(v\mid t)=\vartheta_1(v)$ \quad -- $\vartheta_i$, fonctions de Weierstrass)

% L4–6
Les autres quantités de la théorie des fonctions elliptiques se calculent avec \unsure{avantage|au moyen} à partir des fonctions $\Theta$, une fois fixés $\omega_1$ et $t=\frac{\omega_2}{\omega_1}$.

% L7
Le calcul des $e_i$ par exemple donne

% L8
\[
e_i=\left(\frac1{2\omega_1}\right)^2
\left(\frac13\frac{\Theta'''0}{\Theta'0}-\frac{\Theta_i'''0}{\Theta_i'0}\right)
\quad\text{ou aussi}
\]

% L9–11
\[
\left\{\begin{aligned}
e_1&=\frac13\left(\frac{\pi}{2\omega_1}\right)^2\left[(\Theta_2 0)^4+(\Theta_3 0)^4\right]\\
e_2&=\frac13\left(\frac{\pi}{2\omega_1}\right)^2\left[-(\Theta_1 0)^4-(\Theta_3 0)^4\right]\\
e_3&=\frac12\left(\frac{\pi}{2\omega_1}\right)^2\left[(\Theta_1 0)^4-(\Theta_{\unsure{2|3}}0)^4\right]
\end{aligned}\right.
\tag*{\textcircled{\scriptsize 15}}
\]

% L12
Enfin, on a les relations

% L13
\[
2\eta_1\omega_1=-\frac{\Theta'''0}{6\Theta'0}\quad;\qquad
\Theta'0=\pi\,\Theta_1 0\,\Theta_2 0\,\Theta_3 0
\]

% lower left
% archive sheet number: 89
\scanpage{91}
% L1
\[
\omega_1\quad\omega_2\qquad\left[\frac{\omega_1}{\omega_2}\notin\unsure{R|\mathbb R}\right]
\]

% L2
\[
g_2\quad g_3\qquad\boxed{g_2=20c_1\qquad g_3=28c_2}\qquad(g_2^3-27g_3^2\ne0)
\]

% L3
\[c_1\quad c_2\]

% L4
\[
e_1\quad e_2\quad e_3\qquad 4z^3-g_2z\unsure{-|+}g_3=4(z-e_1)(z-e_2)(z-e_3)
\]

% right of L4
\margin{(\uline{les $e_{\unsure{i|2}}$} \uline{distincts})}
\noindent\rule{.4\linewidth}{.4pt}

% L5
Expressions en fonction de $\omega_1,\omega_2$

% L6–7
\[
\begin{aligned}
c_1&=3\sum_{\omega\in{\unsure{\mathfrak P|\mathscr P}}^{*}}\left(\frac1{2\omega}\right)^4,&
c_2&=5\sum_{\omega\in\mathfrak P^{*}}\left(\frac1{2\omega}\right)^6\\
g_2&=20c_1&g_3&=28c_2
\end{aligned}
\]

% L8
\[
e_i=\unsure{p|\wp}\,\omega_i=\frac1{\omega_i^2}+\sum_{\omega\in\mathfrak P^{*}}\frac1{(\omega_i+2\omega)^2}
\]

% L9
Expressions en fonction de $g_2,g_3$ \deleted{$\unsure{e_3|g_3}$}

% L10
\[
c_1=\frac1{20}g_{\unsure{?|2}},\qquad c_2=\frac1{28}g_3
\]

% L11
$e_1,e_{\unsure{3|2}},e_3$ racines de l'équation $4z^3-g_2z-g_3=0$

% L12
\[
\omega_i=\int_{\infty}^{e_{\unsure{i|2}}}\frac{dz}{\sqrt{4z^3-g_2z-g_3}}\qquad(i,j,k\ \text{distincts})
\]

% L13, under integral
déterminations telles que \deleted{$\sum_{i=1,2,3}\omega_i=0$} \inserted{$\omega_1+\omega_2+\omega_3=0$}

% L14, closely spaced addition
\inserted{\unsure{ou} $\omega_i=\int_{e_j}^{e_k}\frac{dz}{\sqrt{4z^3-g_2z\unsure{-|+}g_3}}$ pour des déterminations convenables.}

% L15
Expressions en \unsure{fonction de} $e_1,e_2,e_3$

% L16
\[
g_2=4(e_2e_3+e_3e_1+e_1e_2)\qquad g_3=4e_1e_2e_3
\]

% L17
le reste alors comme ci-dessus

% L18
Expressions en fonction de $c_1,c_2$

% L19–20
$g_1=20c_1$, $g_{\unsure{2|3}}=20c_{\unsure{3|2}}$, le reste comme ci-dessus.

% lower left
% archive sheet number: 90
\scanpage{92}
% L1, underlined
\uline{Fonction modulaire}

% L2–10
En considérant $g_2,g_3$ comme fonctions de $\omega_1,\omega_2$, où $\unsure{\mathfrak I|J}\frac{\omega_2}{\omega_1}>0$, qui sont homogènes de degré $2$ et $3$ respectivement, on peut former la fonction $\frac{g_2^3}{g_2^3-27g_3^2}$ (finie puisque $\Delta=g_2^3-27g_3^2\ne0$), qui est homogène de degré zéro, donc ne dépend que de $\tau=\frac{\omega_2}{\omega_1}$. On \unsure{l'appelle} fonction modulaire, notation $J(\tau)$; $J(\tau)$ est holomorphe pour $\mathfrak I\tau>0$, et par définition, on a

% L11
\[
\text{pour }\tau=\frac{\omega_2}{\omega_1}\qquad
J(\tau)=\unsure{{}|-}\frac{g_2^3}{\Delta}=\frac1{1-27\frac{g_3^2}{g_2^3}}
\]

% L12–21
Comme $g_2$ et $g_3$ peuvent être pris arbitraires, à \deleted{\unsure{fait|part}} la condition $\Delta=g_2^3-27g_3^2\ne0$ près, $J(\tau)$ prend \deleted{\unsure{toutes}} les valeurs; de plus, le fait que $g_2$ et $g_3$ ne dépendent en fait que du groupe engendré par $\omega_1$ et $\omega_2$ s'exprime par le fait que $J(\tau)$ est invariant par \unsure{les|le} \deleted{\unsure{subst|sous}} groupe arithmétique unimodulaire -- groupe des substitutions $\tau\to\frac{a\unsure{\tau|z}+b}{c\tau+d}$, où $a,b,c,d$ sont entiers, et $ad-bc=1$. De plus,

% lower left
% archive sheet number: 91
\scanpage{93}
% L1–6
il est immédiat que la donnée de $J\tau=a$ ($a$ arbitraire dans $C$) détermine $\tau$ à une transformation arithmétique près, (car il détermine à une homothétie près le groupe des périodes engendré par $\omega_1$ et $\omega_2$, comme on vérifie très facilement).

% L6–11
Il suit aussitôt que $J$ est \deleted{\unsure{univoque}}\inserted{biunivoque} au voisinage de \deleted{tout} point du demi-plan supérieur qui n'est point double d'aucune transformation $\left(\begin{smallmatrix}a,b\\c,d\end{smallmatrix}\right)$ arithmétique, et en ces points seulement.

% L11–14
\deleted{\unsure{Dans ce cas|Donc de plus|Des autres}}\inserted{On} sait d'autre part qu'il y a exactement deux classes de points doubles, savoir $i$ $(i=-\frac1i)$ et $\frac12+i\frac{\sqrt3}{2}$.

% L14–17
Par suite, les seuls points singuliers \inserted{à distance finie} de la fonction inverse $\tau(\unsure{w|u})$ de $J(\tau)$ sont les deux points $J(i)$ et $J(\frac12+i\frac{\sqrt3}{2})$.

% L18–24
Le calcul donne respectivement $0$ et $1$, i.e $g_2=0$ et $g_3=0$ (car en partant de $g_2=0$ resp$^t$ $g_3=0$, on trouve bien $\frac{\omega_2}{\omega_1}=i$ resp$^t$ $\frac{\omega_2}{\omega_1}=\frac1{\unsure{2|3}}+i\frac{\sqrt3}{2}$, par un calcul d'intégrale \deleted{\unsure{a|e}}très simple). -- Ces propriétés sont à la base des théorèmes de Picard sur les fonctions entières.

% lower left
% archive sheet number: 92
\scanpage{94}
% L1, centered title
\begin{center}Fonctions elliptiques de Legendre\end{center}

% L2–3
$\operatorname{sn}u$ est fonction inverse de $\int_0^{\unsure{v|u}}\frac{dx}{\sqrt{(1-x^2)(1-k^2x^2)}}$, impaire. On pose

% L4–6
\[
\left.\begin{aligned}
\operatorname{cn}^2u+\operatorname{sn}^2u&=1\\
\operatorname{dn}^2u+k^2\operatorname{sn}^2u&=1\\
(\operatorname{dn}^2u-k^2\operatorname{cn}^2u&=1-k^2=k'^{\,2})
\end{aligned}\right\}
\quad\begin{aligned}\operatorname{sn}0&=0\\\operatorname{cn}0&=\operatorname{dn}0=1\end{aligned}
\tag*{\textcircled{\scriptsize 1}}
\]

% L7
($\operatorname{cn}$ et $\operatorname{dn}$ sont uniformes, paires)

% L8
On a les équations différentielles

% L9–11
\[
\left\{\begin{aligned}
\operatorname{sn}'u&=\sqrt{(1-\operatorname{sn}^2u)(1-k^2\operatorname{sn}^2u)}=\operatorname{cn}u\operatorname{dn}u\\
\operatorname{cn}'u&=\sqrt{(1-\operatorname{cn}^2u)(k'^{\,2}+k^2\operatorname{cn}^2u)}=-\operatorname{sn}u\operatorname{dn}u\\
\operatorname{dn}'u&=\sqrt{(1-\operatorname{dn}^2u)(\operatorname{dn}^2u-k'^{\,2})}=-k^2\operatorname{sn}u\operatorname{cn}u
\end{aligned}\right.
\tag*{\textcircled{\scriptsize 2}}
\]

% L12
Les développements de Taylor sont, en $0$

% L13–15
\[
\left\{\begin{aligned}
\operatorname{sn}u&=u-\frac{2\alpha k}{3!}u^3+\frac{4k^2(\alpha^2+3)}{5!}u^5-\frac{8k^3(\alpha^3+33\alpha)}{7!}u^7+\ldots\\
\operatorname{cn}u&=1-\frac{u^2}{2!}+\frac{1+4k^2}{4!}u^4-\frac{1+44k^2+16k^4}{6!}u^6+\text{---}\\
\operatorname{dn}u&=1-\frac{k^2u^2}{2!}+\frac{k^2(4+k^2)}{4!}u^4-k^2\frac{16+44k^2+k^4}{6!}u^6+\ldots
\end{aligned}\right.
\tag*{\textcircled{\scriptsize 3}}
\]

% L16
\[\left(\alpha=\frac12\left(k+\frac1k\right)\right)\]

% L17–18
\[
\text{\unsure{Dégénérescences|Dégénérescence}}\quad
\left\{\begin{aligned}
k=0&\quad\operatorname{sn}=\sin&\operatorname{cn}&=\cos,&\operatorname{dn}&=1\\
k=1&\quad\operatorname{sn}=\operatorname{th}&\operatorname{cn}&=\frac1{\operatorname{ch}}&\operatorname{dn}&=\frac1{\operatorname{th}}
\end{aligned}\right.
\]

% L19–20, double underline
\uuline{Périodes et valeurs remarquables}\quad Posons

% L19–20, right of heading
\[
K=\int_0^1\frac{\unsure{du|dx}}{\sqrt{(1-x^2)(1-k^2x^2)}}\qquad
K'=\int_0^1\frac{dx}{\sqrt{(1-x^2)(1-k'^{\,2}x^2)}}
\]

% L21–22, foot
$\operatorname{sn}u,\ldots$ ont les périodes $4K$ et $4iK'$. Posons $K''=-\unsure{k|K}-iK'$, ($4K''$ est une autre période)

% lower left
% archive sheet number: 93
\scanpage{95}
% L1–3, upper left
\[
\operatorname{sn} u\left\{\begin{aligned}
&\text{périodes}&4K&\quad 2iK'\\
&\text{zéros}&2\unsure{m|n}K&+2m'iK'\\
&\text{pôles}&iK'+2mK&+2m'iK'
\end{aligned}\right.
\]

% upper right, beside sn
\margin{si $0\unsure{<|\leq|\deleted{<}}k<1$ \par $\operatorname{sn} u$ n'a pas de zéros ni de pôles pour $u$ réel, et période $4K$}

% L4–6, upper left
\[
\operatorname{cn} u\left\{\begin{aligned}
&\text{périodes}&4K,&-2K''\\
&\text{zéros}&K+2mK&+2m'iK\\
&\text{pôles}&iK'+2mK&+2m'iK'
\end{aligned}\right.
\]

% upper right, beside cn
\margin{pas de pôles pour $u$ réel; zéros $K+2mK$; période $4K$}

% L7–9, upper left
\[
\operatorname{dn} u\left\{\begin{aligned}
&\text{périodes}&2K,&4iK'\\
&\text{zéros}&K''+2mK&+2m'iK\\
&&iK'+2mK&+2m'iK'
\end{aligned}\right.
\]

% upper right, beside dn
\margin{pas de pôles ni de zéros pour $u$ réel; période réelle $2K$}

% L10–11, between blocks
$\operatorname{sn},\operatorname{cn},\operatorname{dn}$ sont, pour leurs périodes propres, d'ordre $2$; pour leurs périodes communes, d'ordre $4$.

% L12–14
\[
\left\{\begin{aligned}
\operatorname{sn}(u+2K)&=-\operatorname{sn} u&\operatorname{cn}(u+2K)&=-\operatorname{cn} u&\operatorname{dn}(u+2K)&=\operatorname{dn} u\\
\operatorname{sn}(u+2iK')&=\operatorname{sn} u&\operatorname{cn}(u+2iK')&=-\operatorname{cn} u&\operatorname{dn}(u+2iK')&=-\operatorname{dn} u\\
\operatorname{sn}(u+2K'')&=-\operatorname{sn} u&\operatorname{cn}(u+2K'')&=\operatorname{cn} u&\operatorname{dn}(u+2K'')&=-\operatorname{dn} u
\end{aligned}\right.
\]

% L15–17
\[
\left\{\begin{aligned}
\operatorname{sn}(u+K)&=\frac{\operatorname{cn} u}{\operatorname{dn} u}&\operatorname{cn}(u+K)&=-k'\frac{\operatorname{sn} u}{\operatorname{dn} u}&\operatorname{dn}(u+K)&=\frac{k'}{\operatorname{dn} u}\\
\operatorname{sn}(u+iK')&=\frac1{k\operatorname{sn} u}&\unsure{\operatorname{cn}|\operatorname{sn}}(u+iK')&=-\frac ik\frac{\operatorname{dn} u}{\operatorname{sn} u}&\operatorname{dn}(u+iK')&=-i\frac{\operatorname{cn} u}{\operatorname{sn} u}\\
\operatorname{sn}(u+\unsure{2K''|K''|\deleted{2}K''})&=\unsure{-|-i}\frac1k\frac{\operatorname{dn} u}{\operatorname{cn} u}&\operatorname{cn}(u+K'')&=-\frac{ik'}{k\operatorname{cn} u}&\operatorname{dn}(u+K'')&=ik'\frac{\operatorname{sn} u}{\operatorname{cn} u}
\end{aligned}\right.
\]

% L18–20
\[
\left\{\begin{aligned}
\operatorname{sn}(iu\mid k)&=i\frac{\operatorname{sn}(u\mid k')}{\operatorname{cn}(u\mid k')}\\
\operatorname{cn}(iu\mid k)&=\frac1{\operatorname{cn}(u\mid k')}\\
\operatorname{dn}(iu\mid k)&=\frac{\operatorname{dn}(u\mid k')}{\operatorname{cn}(u\mid k')}
\end{aligned}\right.
\]

% foot, four rows
\diagram{\scriptsize\setlength{\arraycolsep}{2pt}$
\begin{array}{l|c||c|c|cc||c|c|cc||c|c|cc}
\deleted{\unsure{u|K}}&0&K&2K&3K&\deleted{4K}&iK'&2iK'&3iK'&\deleted{\unsure{4iK'}}&K''&2K''&3K''&\unsure{4|\deleted{4}}\\ \hline
\operatorname{sn} u&0&1&0&-1&&\infty&0&\infty&&-\frac1k&0&\frac1k&\\ \hline
\operatorname{cn} u&1&0&-1&0&&\infty&-1&\infty&&-i\frac{k'}{k}&1&-i\frac{k'}{k}&\\ \hline
\operatorname{dn} u&1&k'&1&k'&&\infty&-1&\infty&&0&-1&0&\\ \hline
\end{array}$}

% lower left
% archive sheet number: 94
\scanpage{96}
% L1–2, centered title
\begin{center}Relations entre les fonctions de Weierstrass et de Legendre\end{center}

% L3–5
Supposons $\omega_1$ réel, $\omega_2$ purement imaginaire ($\unsure{\mathfrak I|J}\omega_2>0$), alors les $e_i$ sont réels; prenons $e_2<e_3<e_1$, posons $k=\sqrt{\frac{e_3-e_2}{e_1-e_2}}$, $M=\sqrt{e_1-e_2}$. Alors

% L6–8
\[
\left\{\begin{aligned}
\operatorname{sn}(u\mid k)&=M\frac{\sigma(\frac uM\mid\omega_1,\omega_2)}{\sigma_2(\frac uM\mid\omega_1,\omega_2)}\\
\operatorname{cn}(u\mid k)&=\deleted{\unsure{M|k}}\frac{\sigma_1(\frac uM\mid\omega_1,\omega_2)}{\sigma_2(\frac uM\mid\omega_1,\omega_2)}\\
\operatorname{dn}(u\mid k)&=\unsure{1|/|,}\,\frac{\sigma_3(\frac uM\mid\omega_1,\omega_2)}{\sigma_1(\frac uM\mid\omega_1,\omega_2)}
\end{aligned}\right.
\]

% L9
\[
\unsure{p|\wp}(u\mid\omega_1,\omega_2)=e_2+\frac{M^2}{\operatorname{sn}^2(Mu\mid k)}
\]

% L10
\[
K=M\omega_1,\quad iK'=M\omega_2,\quad K''=M\omega_3
\]

% L11–16
(Pour $k$ donné, pour \unsure{rattacher} $\operatorname{sn}(u\mid k)\ldots$ aux fonctions de Weierstrass, on prend $M$ arbitraire, puis on détermine $\omega_1,\omega_2,\omega_3$ par $K=M\omega_1$, $iK'=M\omega_{\unsure{2|1}}$, $K''=M\omega_3$, et on a les formules ci-dessus; \deleted{\unsure{où|ici}} alors $e_1,e_2,e_3$ sont donnés par

% L17
\[
3e_1=(1+k'^{\,2})M^2\qquad
3e_2=(\unsure{-2+|2+}k'^{\,2})M^2\qquad
3e_3=(1-2k'^{\,2})M^2
\]

% L18–21, foot
On pourra faire par exemple $M=\frac12$, donc prendre des périodes \inserted{communes} $2\omega_1=4K,\unsure{,|{}}$ $2\omega_2=4iK'$; alors $\operatorname{sn},\operatorname{cn},\operatorname{dn}$ sont fonctions elliptiques d'ordre $4$ pour les périodes $2\omega_1,2\omega_2$)

% lower left
% archive sheet number: 95
\scanpage{97}
% L1, heading
Fonctions de Legendre et fonctions $\Theta$

% L2–4
\[
\left\{\begin{aligned}
\operatorname{sn}(u\mid k)&=\frac1{\sqrt k}\frac{\Theta(\frac{u}{2K}\mid\frac{iK'}K)}{\Theta_2(\frac{u}{2K}\mid\frac{iK'}K)}\\
\operatorname{cn}(u\mid k)&=\sqrt{\frac{k'}k}\frac{\Theta_1(\frac{u}{2K},\frac{iK'}K)}{\Theta_2(\frac{u}{2K},\frac{iK'}K)}\\
\operatorname{dn}(u\mid k)&=\sqrt{k'}\frac{\Theta_3(\frac{u}{2K},\frac{iK'}K)}{\Theta_2(\frac{u}{2K},\frac{iK'}K)}
\end{aligned}\right.
\]

% L5, below two rules
\begin{center}\rule{.5\linewidth}{.4pt}\\[-.5ex]\rule{.27\linewidth}{.4pt}\end{center}
\uline{Formules d'addition}

% L6
\[
\operatorname{sn}(u+v)=\frac{\operatorname{sn} u\,\operatorname{cn} v\,\operatorname{dn} v+\operatorname{sn} v\,\operatorname{cn} u\,\operatorname{dn} u}{1-k^2{\unsure{\mathrm{on}|\operatorname{sn}}}^2u\,\operatorname{sn}^2v}
\]

% L7
\[
\operatorname{cn}(u+v)=\frac{\operatorname{cn} u\,\operatorname{cn} v-\operatorname{sn} u\,\operatorname{sn} v\,\operatorname{dn} u\,\operatorname{dn} v}{1-k^2\operatorname{sn}^2u\,\operatorname{sn}^2v}
\]

% L8
\[
\operatorname{dn}(u+v)=\frac{\operatorname{dn} u\,\operatorname{dn} v-k^2\operatorname{sn} u\,\operatorname{sn} v\,\operatorname{cn} u\,\operatorname{cn} v}{1-k^2\operatorname{sn}^2u\,\operatorname{sn}^2v}
\]

% lower left
% archive sheet number: 96
